\chapter*{全书导读}
\addcontentsline{toc}{chapter}{全书导读}
\markboth{全书导读}{全书导读}

\noindent\textbf{一句话主旨。}全书讲的是黎曼面上局部系统（local system）模空间的两副面孔：
一副是「经典」的——基本群的表示簇、Higgs 丛与 Hitchin 系统；一副是「量子」的——圈群、
可积系统、拓扑场论与几何量子化。四章各取一面，而第 2 章给出的 $\pi_1^{\R}$ 及其表示簇，
正是另三章反复用到的地基。

\section*{第 1 章\quad 圈群、cluster、二聚体与可积系统}

作者按「两种构造、一个同构」来组织：先回顾仿射 Poisson--Lie 群（$\hat A_{N-1}$ 型，
即 $PGL(N)$ 的圈群）上的标准可积系统，再回顾 Goncharov--Kenyon 用二维环面上的二部
「二聚体」（dimer）模型造出的 cluster 可积系统，然后证明两者同构。中段的
\S1.3--\S1.4 用 double Bruhat 胞腔的 cluster 参数化把简单群与圈群统一处理，关键工具是
Cartan--Weyl 生成元、coextended 仿射 Weyl 群与 Thurston 图；\S1.5 把 dimer 分割函数
与离散 Dirac 算子联系起来，得到谱子流形与 Poisson 括号；\S1.8 的例子里有最简单相对论
Toda 链、rank two 相对论 Toda 链，以及五角星映射（pentagram map）。六个附录补足了
cluster 簇、Thurston 图与 Schwartz 坐标的技术细节。

\emph{读法提示：}本章公式密集、图多，宜从 \S1.1 的三小节读起，先抓住「同一个可积系统
的两种坐标」这一主线，再进 \S1.5 的 dimer 机器；附录不必顺读，用到再查。

\section*{第 2 章\quad 克莱因曲面及其基本群（本卷全译）}

三节：\S2.1 讲 Klein 曲面的一般理论（实代数曲线、拓扑分类、双解析结构），\S2.2 定义
克莱因曲面的基本群 $\pi_1^{\R}((X,\sigma);x)$ 并证明它分类实覆盖，\S2.3 讨论线性表示、
酉表示与酉表示簇的实结构。核心是一条短正合列
\[
1\longrightarrow \pi_1(X;x)\longrightarrow \pi_1^{\R}((X,\sigma);x)
\longrightarrow \operatorname{Gal}(\C/\R)\longrightarrow 1,
\]
以及一个类比：SGA-1 意义下代数基本群之于拓扑基本群，正如 $\pi_1^{\R}$ 之于 $\pi_1$。
每一节都配了习题（2.1.17--2.1.27、2.2.22--2.2.31、2.3.5--2.3.10），其中不少是正文结论的
另一证明或具体算例，值得动手。

\emph{读法提示：}若只关心「实结构如何作用在表示簇上」，可直接读 \S2.2.3 与 \S2.3.2，
再回补 \S2.2.2 的构造；定理 2.1.8（Weichold）是全章唯一的分类定理，\S2.1.2 的
$(g,k,a)$ 三元组要记牢。

\section*{第 3 章\quad 拓扑场论五讲}

从 Atiyah 的 TQFT 公理出发（\S3.1），经有限群规范理论、Frobenius 代数、二维
Yang--Mills、上同调场论等例子，过渡到扩展 TQFT 与高阶范畴（\S3.3），再在一维、二维
情形陈述并「看图」证明配边猜想（\S3.4，含对偶对象、Serre 扭、有向与 $r$-spin 理论、
有向把柄的伴随），最后走向高维（\S3.5，含 Drinfeld 中心、融合范畴与 Serre 自同构）。
全章的主线是：一个 $n$ 维扩展 TQFT 的「值」由一点处的对象加上可对偶性决定。

\emph{读法提示：}这是 Lurie 论文的 companion，读法也应如是——先把 \S3.4 的二维情形
吃透（Frobenius 代数与配边猜想的对应最干净），再看 \S3.5 如何把同一套语言升到高维。

\section*{第 4 章\quad 黎曼面上的 Higgs 丛与局部系统}

\S4.2 从 Dolbeault 面讲 Higgs 丛与模空间，并给出 Hitchin--Kobayashi 对应（稳定
Higgs 丛 $\leftrightarrow$ 规范场方程解）；\S4.3 从 Betti 面讲表示簇、局部系统与全纯
联络，给出 Corlette--Donaldson 定理与 hyperkähler 约化，于是得到两个模空间之间的
对应（非阿贝尔 Hodge 对应）；\S4.4 把对应用到微分方程：单值化、高阶方程、opers 与
Eichler--Shimura 同构，收束于几何 Langlands 的入口。

\emph{读法提示：}若已熟悉 Hitchin 系统，可从 \S4.3.3 切入；若想看清「对应」的解析
机制，则 \S4.2.3 与 \S4.3.4 是两根支柱，值得逐行核。

\section*{四章之间的接口}

\begin{itemize}[leftmargin=1.8em, itemsep=0.25em]
\item 第 2 章的 $\pi_1^{\R}$ 与酉表示簇的实结构，是第 4 章「实 Higgs 丛 / 实局部系统」
      的语言基础；第 2 章定理 2.3.4（$\Phi$ 的像落在 $\Fix(\kappa)$）正是「实点」这一
      观念的拓扑原型。
\item 第 1 章的 dimer/cluster 机器与第 3 章的二维 TQFT 共享「二维环面 + 二部图 +
      分割函数」这套组合学，\S1.5.6 的谱子流形与 \S3.3 的融合范畴值得对照。
\item 第 3 章的扩展 TQFT 与第 4 章的几何 Langlands 在「oper 模空间 = 某些 TQFT 的
      状态空间」这一点上相遇，这是当前研究的前沿接口之一。
\end{itemize}

\begin{annote}
本节只作导航：第 1、3、4 章的内容概述据原书前言与目录写成，本卷未收其译文，
引用时请以原书页码为准（第 1 章 p.1--65，第 3 章 p.109--163，第 4 章 p.165--215）。
\end{annote}
